\documentclass[a4paper,12pt]{article}

\usepackage[T2A]{fontenc}
\usepackage[utf8]{inputenc}
\usepackage[russian]{babel}

\usepackage{amsmath,amssymb,amsthm,mathtools}
\usepackage{helvet}
\renewcommand{\familydefault}{\sfdefault}

\usepackage{icomma}
\usepackage{geometry}
\geometry{
  left=20mm,
  right=20mm,
  top=20mm,
  bottom=22mm
}

\usepackage{microtype}
\usepackage{tikz}
\usetikzlibrary{calc}
\usetikzlibrary{arrows.meta,positioning,shapes.geometric}

\usepackage{tkz-euclide}
\usepackage{pgfplots}
\pgfplotsset{compat=1.18}

\usepackage{tabularx,booktabs,longtable}
\usepackage{enumitem}
\usepackage{hyperref}
\usepackage{parskip}
\usepackage{fancyhdr}
\usepackage{setspace}
\usepackage{xcolor}

% ------------------------------------------------
% Цветовая палитра
% ------------------------------------------------

\definecolor{accent}{HTML}{008080}
\definecolor{darkaccent}{HTML}{004D4D}
\definecolor{textgray}{HTML}{333333}
\definecolor{softgray}{HTML}{F2F5F5}

\hypersetup{
  colorlinks=true,
  linkcolor=darkaccent,
  urlcolor=darkaccent
}

\color{textgray}
\onehalfspacing
\setlength{\parindent}{0pt}
\setlength{\emergencystretch}{2em}

\setlist[enumerate]{
  itemsep=0.55em,
  topsep=0.5em,
  leftmargin=*
}

\setlist[itemize]{
  itemsep=0.35em,
  topsep=0.4em,
  leftmargin=*
}

\renewcommand{\arraystretch}{1.35}

% ------------------------------------------------
% Служебные команды
% ------------------------------------------------

\newcommand{\Prob}{\mathsf{P}}

\newcommand{\topichead}[1]{%
  \vspace{0.5em}
  {\Large\bfseries\color{darkaccent}#1\par}
  \vspace{0.35em}
  {\color{accent}\rule{\linewidth}{0.7pt}}
  \vspace{0.5em}
}

\newcommand{\subtopic}[1]{%
  \vspace{0.4em}
  {\large\bfseries\color{darkaccent}#1\par}
  \vspace{0.15em}
}

% ------------------------------------------------
% Стили блок-схем
% ------------------------------------------------

\tikzset{
  flowstart/.style={
    ellipse,
    draw=darkaccent,
    line width=0.9pt,
    minimum width=4.8cm,
    minimum height=1.1cm,
    align=center,
    font=\small\bfseries
  },
  flowproc/.style={
    rectangle,
    rounded corners=3pt,
    draw=darkaccent,
    line width=0.9pt,
    minimum width=5.4cm,
    minimum height=1.15cm,
    text width=5.4cm,
    inner sep=5pt,
    align=center,
    font=\small
  },
  flowcond/.style={
    diamond,
    aspect=2.15,
    draw=darkaccent,
    line width=0.9pt,
    text width=5.1cm,
    inner sep=1.5pt,
    align=center,
    font=\small
  },
  flowarrow/.style={
    -{Latex[length=3mm,width=2mm]},
    line width=0.9pt,
    draw=darkaccent
  },
  flowlabel/.style={
    font=\small\bfseries,
    text=darkaccent,
    fill=white,
    inner sep=1.5pt
  }
}

% ------------------------------------------------
% Колонтитулы
% ------------------------------------------------

\pagestyle{fancy}
\fancyhf{}
\fancyhead[L]{\small Ксения Клыкова}
\fancyhead[R]{\small ЕГЭ $\cdot$ теория вероятностей}
\fancyfoot[C]{\small\thepage}
\renewcommand{\headrulewidth}{0.4pt}
\renewcommand{\footrulewidth}{0pt}

% ================================================================
\begin{document}

% ================================================================
% Титульный лист
% ================================================================

\begin{titlepage}
\thispagestyle{empty}

\begin{center}

\vspace*{2.8cm}

{\large МАТЕМАТИКА\par}

\vspace{1.2cm}

{\fontsize{28}{34}\selectfont
\bfseries\color{darkaccent}
Чек-лист\par}

\vspace{0.8cm}

{\LARGE\bfseries
Теория вероятностей\par}

\vspace{0.4cm}

{\Large
дерево вероятностей, повторные испытания,\\
полная вероятность и противоположные события\par}

\vspace{1.8cm}

{\large
Подготовка к ЕГЭ по профильной математике\par}

\vfill

{\large
Ксения Клыкова\par}

\vspace{0.8cm}

{\normalsize
Лёвин Артём Александрович\\
эксклюзивно для Ксении Клыковой\par}

\vspace{0.8cm}

{\normalsize\today\par}

\vspace{2.1cm}

\begin{tikzpicture}[remember picture,overlay]
  \draw[
    accent,
    line width=1.2pt
  ]
  ([xshift=20mm,yshift=24mm]current page.south west)
  .. controls
  ([xshift=55mm,yshift=36mm]current page.south west)
  and
  ([xshift=-55mm,yshift=12mm]current page.south east)
  ..
  ([xshift=-20mm,yshift=24mm]current page.south east);
\end{tikzpicture}

\end{center}
\end{titlepage}

% ================================================================
% Основной чек-лист
% ================================================================

\topichead{1. Язык вероятностной задачи}

Большинство задач на дерево вероятностей сводятся к правильному переводу
словесного условия на язык событий.

\begin{center}
\begin{tabularx}{\linewidth}{@{}>{\bfseries}p{0.18\linewidth}X X@{}}
\toprule
Фраза & Математический смысл & Что обычно делаем \\
\midrule
«и» &
события происходят последовательно &
умножаем вероятности вдоль одного маршрута \\

«или» &
подходит один из нескольких несовместимых маршрутов &
складываем вероятности маршрутов \\

«не» &
противоположное событие &
используем $1-\Prob(A)$ \\

«хотя бы один» &
один, два, три и т.\,д. успешных исхода &
часто выгодно считать через противоположное событие «ни одного» \\

«не менее» &
вероятность должна достичь заданного порога &
составляем неравенство и ищем минимальное число испытаний \\
\bottomrule
\end{tabularx}
\end{center}

\subtopic{Главное правило маршрута}

Для последовательных событий действует формула
\[
\Prob(A\cap B)=\Prob(A)\Prob(B\mid A).
\]

Поэтому вероятности вдоль одной ветви дерева \textbf{перемножаются}.

Если события независимы, то
\[
\Prob(B\mid A)=\Prob(B),
\]
и формула принимает привычный вид
\[
\Prob(A\cap B)=\Prob(A)\Prob(B).
\]

\subtopic{Когда вероятности складываются}

Если событие может произойти по нескольким \textbf{несовместимым}
маршрутам $R_1,R_2,\ldots,R_k$, то
\[
\Prob(A)
=
\Prob(R_1)+\Prob(R_2)+\cdots+\Prob(R_k).
\]

Различные законченные маршруты дерева обычно несовместимы: одновременно
реализоваться могут только события одного маршрута.

\subtopic{Типичная ошибка}

Фраза «или» сама по себе ещё не гарантирует простое сложение.

В общем случае
\[
\Prob(A\cup B)
=
\Prob(A)+\Prob(B)-\Prob(A\cap B).
\]

Просто складывать $\Prob(A)$ и $\Prob(B)$ можно, если события
$A$ и $B$ несовместимы.

% ================================================================

\topichead{2. Дерево вероятностей}

Дерево удобно использовать, когда эксперимент проходит в несколько этапов:
выстрелы, проверки, игры, выбор устройства, выбор завода, повторный анализ.

\subtopic{Правила построения}

\begin{enumerate}[label=\arabic*)]
\item Каждому этапу соответствует новый уровень дерева.
\item Из одной вершины выходят все возможные исходы следующего этапа.
\item Сумма вероятностей ветвей, выходящих из одной вершины, равна $1$.
\item Вероятность полного маршрута равна произведению вероятностей его ветвей.
\item Если событию соответствуют несколько маршрутов, вероятности этих
маршрутов складывают.
\end{enumerate}

\subtopic{Мини-пример}

Если вероятность успеха равна $0{,}6$, то вероятность неудачи равна
\[
1-0{,}6=0{,}4.
\]

Если после первой неудачи проводится ещё одна попытка с вероятностью успеха
$0{,}8$, то вероятность маршрута

\[
\text{«неудача на первой попытке и успех на второй»}
\]

равна
\[
0{,}4\cdot0{,}8=0{,}32.
\]

Вероятность успеха не позднее второй попытки:
\[
0{,}6+0{,}4\cdot0{,}8
=
0{,}92.
\]



% ================================================================

\topichead{3. Повторные попытки до первого успеха}

В задачах с повторными выстрелами или проверками важно различать:

\begin{itemize}
\item вероятность успеха \textbf{ровно на $k$-й попытке};
\item вероятность успеха \textbf{не позднее $k$-й попытки}.
\end{itemize}

Это разные события.

\subtopic{Успех ровно на второй попытке}

Если
\[
p_1=0{,}6,
\qquad
p_2=0{,}8,
\]
то успех ровно на второй попытке означает:

\[
\text{первая попытка неудачна, вторая успешна}.
\]

Следовательно,
\[
\Prob(X=2)
=
(1-0{,}6)\cdot0{,}8
=
0{,}32.
\]

\subtopic{Успех не позднее второй попытки}

Подходят два несовместимых маршрута:

\[
\text{успех на первой}
\qquad\text{или}\qquad
\text{неудача на первой и успех на второй}.
\]

Поэтому
\[
\Prob(X\le2)
=
0{,}6+0{,}4\cdot0{,}8
=
0{,}92.
\]

\subtopic{Успех не позднее третьей попытки}

\[
\Prob(X\le3)
=
0{,}6
+
0{,}4\cdot0{,}8
+
0{,}4\cdot0{,}2\cdot0{,}8.
\]

Отсюда
\[
\Prob(X\le3)
=
0{,}984.
\]

Для четырёх попыток:
\[
\Prob(X\le4)
=
0{,}984
+
0{,}4\cdot0{,}2\cdot0{,}2\cdot0{,}8
=
0{,}9968.
\]

Следовательно, чтобы вероятность успеха была не менее $0{,}99$,
потребуется минимум
\[
\boxed{4}
\]
попытки.

% ================================================================

\topichead{4. Быстрый способ: противоположное событие}

Предыдущий расчёт можно существенно сократить.

Событие

\[
A=\{\text{успех произошёл хотя бы один раз}\}
\]

противоположно событию

\[
\overline A=\{\text{все попытки оказались неудачными}\}.
\]

Поэтому
\[
\Prob(A)=1-\Prob(\overline A).
\]

\subtopic{Одинаковая вероятность успеха}

Если каждая попытка независима и вероятность успеха всегда равна $p$, то
\[
\Prob(\text{ни одного успеха за }n\text{ попыток})
=
(1-p)^n.
\]

Следовательно,
\[
\boxed{
\Prob(\text{хотя бы один успех})
=
1-(1-p)^n
}.
\]

\subtopic{Пример}

При $p=0{,}8$ требуется достичь вероятности не менее $0{,}95$.

За одну попытку:
\[
0{,}8<0{,}95.
\]

За две попытки:
\[
1-(1-0{,}8)^2
=
1-0{,}2^2
=
0{,}96.
\]

Следовательно,
\[
\boxed{n=2}.
\]

\subtopic{Первая вероятность отличается от последующих}

Пусть первая попытка имеет вероятность успеха $p_1$, а каждая следующая ---
$p$.

Вероятность всех неудач за $n$ попыток:
\[
(1-p_1)(1-p)^{n-1}.
\]

Поэтому
\[
\boxed{
\Prob(\text{успех не позднее }n\text{-й попытки})
=
1-(1-p_1)(1-p)^{n-1}
}.
\]

Для
\[
p_1=0{,}6,
\qquad
p=0{,}8
\]
получаем
\[
\Prob_n
=
1-0{,}4\cdot0{,}2^{\,n-1}.
\]

При $n=3$:
\[
\Prob_3
=
1-0{,}4\cdot0{,}2^2
=
0{,}984.
\]

При $n=4$:
\[
\Prob_4
=
1-0{,}4\cdot0{,}2^3
=
0{,}9968.
\]

% ================================================================
% Блок-схема 2
% ================================================================

\clearpage
\thispagestyle{fancy}

\begin{center}
{\Large\bfseries\color{darkaccent}
Алгоритм: минимальное число попыток для достижения заданной вероятности}
\end{center}

\vspace{10mm}

\begin{center}
\begin{tikzpicture}[node distance=10mm]

\node[flowstart] (start)
{Начало};

\node[flowproc, below=of start] (formula)
{Записать вероятность полного отсутствия успеха после $n$ попыток};

\node[flowproc, below=of formula] (target)
{Записать
$\Prob_n=1-\Prob(\text{все попытки неудачны})$};

\node[flowproc, below=of target] (setn)
{Начать с минимального допустимого значения $n$};

\node[flowcond, below=11mm of setn] (check)
{Выполнено ли условие $\Prob_n\ge p_{\text{треб}}$?};

\node[flowproc, right=32mm of check] (increase)
{Увеличить $n$ на $1$ и пересчитать вероятность};

\node[flowproc, below=16mm of check] (answer)
{Зафиксировать найденное минимальное значение $n$};

\node[flowstart, below=of answer] (finish)
{Ответ};

\draw[flowarrow] (start) -- (formula);
\draw[flowarrow] (formula) -- (target);
\draw[flowarrow] (target) -- (setn);
\draw[flowarrow] (setn) -- (check);

\draw[flowarrow]
(check.south) --
node[flowlabel,right]{да}
(answer.north);

\draw[flowarrow]
(check.east) --
node[flowlabel,above]{нет}
(increase.west);

\draw[flowarrow]
(increase.north) |- (target.east);

\draw[flowarrow] (answer) -- (finish);

\end{tikzpicture}
\end{center}

\clearpage

% ================================================================

\topichead{5. Формула полной вероятности}

Если событие $B$ может произойти после одного из нескольких взаимоисключающих
сценариев
\[
A_1,A_2,\ldots,A_n,
\]
образующих полную группу событий, то
\[
\boxed{
\Prob(B)
=
\sum_{i=1}^{n}
\Prob(A_i)\Prob(B\mid A_i)
}.
\]

Для двух ветвей:
\[
\Prob(B)
=
\Prob(A_1)\Prob(B\mid A_1)
+
\Prob(A_2)\Prob(B\mid A_2).
\]

\subtopic{Пример с двумя ветвями}

Пусть
\[
\Prob(A_1)=0{,}4,
\qquad
\Prob(A_2)=0{,}6,
\]
а условные вероятности события $B$ равны
\[
\Prob(B\mid A_1)=0{,}1,
\qquad
\Prob(B\mid A_2)=0{,}8.
\]

Тогда
\[
\Prob(B)
=
0{,}4\cdot0{,}1
+
0{,}6\cdot0{,}8
=
0{,}52.
\]

Здесь сначала определяется, по какой ветви пошёл эксперимент, затем ---
произошло ли событие $B$ внутри выбранной ветви.

% ================================================================

\topichead{6. Неизвестная вероятность первой ветви}

Иногда вероятность выбора одного из вариантов в условии отсутствует.

Если вариантов ровно два, удобно ввести
\[
\Prob(A_1)=x,
\qquad
\Prob(A_2)=1-x.
\]

После этого формула полной вероятности превращается в обычное уравнение.

\subtopic{Пример: два завода}

Пусть лампа произведена на первом заводе с вероятностью $x$, на втором ---
с вероятностью $1-x$.

Доля брака:

\[
2\% \quad\text{на первом заводе},
\qquad
3\% \quad\text{на втором}.
\]

Общая вероятность получить бракованную лампу равна
\[
0{,}024.
\]

Получаем уравнение
\[
0{,}02x+0{,}03(1-x)=0{,}024.
\]

Раскрываем скобки:
\[
0{,}02x+0{,}03-0{,}03x=0{,}024.
\]

Отсюда
\[
-0{,}01x=-0{,}006,
\qquad
\boxed{x=0{,}6}.
\]

Значит, вероятность того, что случайно выбранная лампа произведена первым
заводом, равна $0{,}6$.

\subtopic{Что обязательно проверить}

Для вероятности $x$ должно выполняться
\[
0\le x\le1.
\]

Если после решения уравнения получено число вне этого промежутка, значит,
в модели или вычислениях допущена ошибка.

% ================================================================

\topichead{7. Две игры: дерево или таблица}

Если на каждом этапе имеется несколько одинаковых по смыслу исходов,
таблица иногда компактнее дерева.

Пусть в каждой из двух независимых игр команда:

\[
\begin{array}{ccl}
\text{выигрывает} &:& 0{,}45,\\
\text{играет вничью} &:& 0{,}10,\\
\text{проигрывает} &:& 0{,}45.
\end{array}
\]

За победу начисляется $3$ очка, за ничью --- $1$, за поражение --- $0$.

Команда проходит дальше, если набирает не менее $4$ очков.

\subtopic{Удобный расчёт через противоположное событие}

Получить не менее $4$ очков можно только тремя способами:

\[
(3;3),\qquad(3;1),\qquad(1;3).
\]

Следовательно,
\[
\Prob(\text{пройти})
=
0{,}45^2
+
0{,}45\cdot0{,}10
+
0{,}10\cdot0{,}45.
\]

Получаем
\[
\Prob(\text{пройти})
=
0{,}2025+0{,}045+0{,}045
=
0{,}2925.
\]

Поэтому
\[
\boxed{
\Prob(\text{не пройти})
=
1-0{,}2925
=
0{,}7075
}.
\]

\subtopic{Почему таблица удобна}

Для двух игр имеется всего
\[
3\cdot3=9
\]
комбинаций результатов.

В таблице их можно увидеть сразу. Дерево даёт ту же математическую модель,
однако занимает больше места.

Основная идея сохраняется:

\[
\boxed{
\text{вдоль маршрута --- умножение,\qquad
между подходящими маршрутами --- сложение}.
}
\]

% ================================================================

\topichead{8. Повторная проверка только после отрицательного результата}

Особенно внимательно нужно читать условие, если второй этап проводится
только после определённого исхода первого.

Пусть анализ даёт положительный результат с вероятностью
\[
0{,}4,
\]
а отрицательный --- с вероятностью
\[
0{,}6.
\]

После отрицательного результата анализ повторяют один раз с теми же
вероятностями.

Выявление произойдёт либо:

\[
\text{сразу на первом анализе},
\]

либо

\[
\text{после первого отрицательного результата на втором анализе}.
\]

Поэтому
\[
\Prob(\text{выявить})
=
0{,}4
+
0{,}6\cdot0{,}4.
\]

Следовательно,
\[
\boxed{
\Prob(\text{выявить})=0{,}64
}.
\]

Ветвь повторной проверки строится только от результата, после которого
повтор действительно назначается.

% ================================================================

\topichead{9. «Хотя бы один»: главный приём}

Формулировка «хотя бы один» означает:

\[
1,\ 2,\ 3,\ldots
\]

подходящих объектов или успешных исходов.

Вместо перечисления всех вариантов обычно проще использовать
противоположное событие.

\subtopic{Пример с двумя лампами}

Вероятность перегорания каждой лампы в течение года равна
\[
0{,}3.
\]

При независимости ламп вероятность того, что конкретная лампа
\textbf{не перегорит}, равна
\[
0{,}7.
\]

Вероятность того, что не перегорит ни одна из двух ламп:
\[
0{,}7^2=0{,}49.
\]

Следовательно,
\[
\boxed{
\Prob(\text{хотя бы одна перегорит})
=
1-0{,}49
=
0{,}51
}.
\]

\subtopic{Важная ловушка}

Вычисление
\[
1-0{,}3^2=0{,}91
\]
здесь отвечает другому вопросу:

\[
\text{«какова вероятность, что перегорят не обе лампы?»}
\]

Это событие не совпадает с событием
«хотя бы одна лампа перегорит».

% ================================================================

\topichead{10. Как выбирать между деревом, таблицей и дополнением}

\begin{center}
\begin{tabularx}{\linewidth}{@{}p{0.25\linewidth}X@{}}
\toprule
Ситуация & Удобный инструмент \\
\midrule

Несколько последовательных этапов &
дерево вероятностей \\

На каждом из двух этапов несколько однотипных исходов &
таблица или дерево \\

Нужно «хотя бы одно», «не менее одного» &
противоположное событие \\

Нужно добиться вероятности не ниже заданной &
формула вероятности неудачи и поиск минимального $n$ \\

Есть несколько источников или способов выбора &
формула полной вероятности \\

Вероятность одной из двух первых ветвей неизвестна &
обозначить её через $x$, вторую --- через $1-x$ \\

Количество неудобных исходов существенно больше количества удобных &
посчитать противоположное событие \\
\bottomrule
\end{tabularx}
\end{center}

% ================================================================

\topichead{11. Типичные ошибки}

\begin{enumerate}[label=\arabic*)]
\item Складывать вероятности событий, которые должны происходить
последовательно.

\item Умножать вероятности разных альтернативных маршрутов вместо их сложения.

\item Считать только последний успешный выстрел и забывать, что до него
необходимо пройти через предыдущие неудачи.

\item Путать
\[
\text{«успех на третьей попытке»}
\]
и
\[
\text{«успех не позднее третьей попытки»}.
\]

\item Забывать, что вероятности ветвей, выходящих из одной вершины дерева,
должны давать в сумме $1$.

\item Использовать формулу
\[
1-p^n
\]
для события «хотя бы один успех».
При независимых одинаковых испытаниях правильная формула:
\[
1-(1-p)^n.
\]

\item При вводе неизвестной вероятности $x$ забывать, что для двух
исчерпывающих вариантов вторая вероятность равна $1-x$.

\item После решения задачи на минимальное число попыток записывать в ответ
полученную вероятность вместо количества попыток.

\item Использовать независимость событий без основания.
Если вероятность второго события зависит от первого, следует работать с
условной вероятностью
\[
\Prob(B\mid A).
\]

\item Не проверять итог:
\[
0\le\Prob(A)\le1.
\]
\end{enumerate}

% ================================================================

\topichead{12. Финальная памятка}

Перед записью ответа проверьте:

\begin{itemize}
\item Что является отдельным этапом эксперимента?
\item Какие исходы возможны на каждом этапе?
\item Суммируются ли вероятности ветвей из одной вершины до $1$?
\item Какие маршруты соответствуют нужному событию?
\item Где в формулировке скрыто «и»?
\item Где скрыто «или»?
\item Можно ли быстрее посчитать противоположное событие?
\item Если требуется минимальное число испытаний, проверено ли предыдущее
значение $n$?
\item Соответствует ли записанный ответ тому, что спрашивается в условии?
\end{itemize}

Ключевая схема темы:
\[
\boxed{
\begin{gathered}
\text{ветви одного маршрута}\ \Longrightarrow\ \text{умножаем},\\[2mm]
\text{несовместимые подходящие маршруты}\ \Longrightarrow\ \text{складываем},\\[2mm]
\text{«хотя бы один»}\ \Longrightarrow\
1-\Prob(\text{ни одного}).
\end{gathered}
}
\]

% ================================================================
% Тренировочный блок
% ================================================================

\clearpage

\begin{center}
{\LARGE\bfseries\color{darkaccent}
Тренировочный блок}
\par\vspace{0.3em}
{\large Путь А}
\end{center}

\vspace{0.7em}

Задачи расположены по нарастающей сложности. Решения в тренировочном
блоке не приводятся.

\subtopic{Средний уровень}

\begin{enumerate}[label=\arabic*)]

\item
Вероятность успешного выполнения одной попытки равна $0{,}7$.
Проводятся две независимые попытки.
Найдите вероятность того, что хотя бы одна попытка окажется успешной.

\item
Вероятность успеха при первой попытке равна $0{,}5$, а при каждой
следующей --- $0{,}8$.
Попытки продолжаются до первого успеха.
Найдите вероятность того, что успех произойдёт не позднее третьей попытки.

\item
В помещении установлены две независимые лампы.
Вероятность перегорания каждой лампы в течение года равна $0{,}2$.
Найдите вероятность того, что за год перегорит хотя бы одна лампа.

\end{enumerate}

\subtopic{Выше среднего}

\begin{enumerate}[label=\arabic*),resume]

\item
Вероятность попадания при каждом независимом выстреле равна $0{,}75$.
Какое минимальное число выстрелов необходимо сделать, чтобы вероятность
хотя бы одного попадания была не менее $0{,}99$?

\item
Вероятность успеха при первой попытке равна $0{,}4$, а при каждой
следующей --- $0{,}7$.
Найдите минимальное число попыток, при котором вероятность получить успех
хотя бы один раз составляет не менее $0{,}98$.

\item
Изделие поступает с первой производственной линии с вероятностью $0{,}35$,
а со второй --- с вероятностью $0{,}65$.
Вероятность брака на первой линии равна $0{,}02$, на второй --- $0{,}05$.
Найдите вероятность того, что случайно выбранное изделие окажется
бракованным.

\item
Некоторое изделие изготовлено на первом заводе с вероятностью $x$, а на
втором --- с вероятностью $1-x$.
Вероятность брака на первом заводе равна $0{,}01$, на втором --- $0{,}04$.
Общая вероятность брака равна $0{,}022$.
Найдите $x$.

\item
Футбольная команда играет две независимые игры.
В каждой игре вероятность победы равна $0{,}5$, ничьей --- $0{,}2$,
поражения --- $0{,}3$.
За победу начисляют $3$ очка, за ничью --- $1$, за поражение --- $0$.
Найдите вероятность того, что за две игры команда наберёт меньше $4$
очков.

\item
Вероятность положительного результата одного анализа равна $0{,}35$.
Если результат отрицательный, анализ повторяют; всего разрешено не более
трёх анализов.
Считая результаты независимыми, найдите вероятность получить хотя бы один
положительный результат.

\end{enumerate}

\subtopic{Сложная задача}

\begin{enumerate}[label=\arabic*),resume]

\item
Стрелок случайно выбирает одно из двух устройств и после выбора делает из
него два независимых выстрела.

Первое устройство выбирается с вероятностью $0{,}4$, второе --- с
вероятностью $0{,}6$.

Вероятность попадания одним выстрелом из первого устройства равна
$0{,}9$, а из второго --- $0{,}5$.

Найдите вероятность того, что хотя бы один из двух выстрелов попадёт в
цель.

\end{enumerate}

% ================================================================
% Ответы
% ================================================================

\clearpage

\begin{center}
{\LARGE\bfseries\color{darkaccent}
Ответы к тренировочному блоку}
\end{center}

\vspace{1em}

\begin{center}
\begin{tabularx}{0.88\linewidth}{
@{}
>{\centering\arraybackslash}p{0.11\linewidth}
>{\centering\arraybackslash}X
>{\centering\arraybackslash}p{0.11\linewidth}
>{\centering\arraybackslash}X
@{}
}
\toprule
\textbf{№} & \textbf{Ответ} & \textbf{№} & \textbf{Ответ} \\
\midrule
1 & $0{,}91$     & 6  & $0{,}0395$   \\
2 & $0{,}98$     & 7  & $0{,}6$      \\
3 & $0{,}36$     & 8  & $0{,}55$     \\
4 & $4$          & 9  & $0{,}725375$ \\
5 & $4$          & 10 & $0{,}846$    \\
\bottomrule
\end{tabularx}
\end{center}

\vfill

\begin{center}
{\small\color{textgray}
Лёвин Артём Александрович
\quad$\cdot$\quad
Ксения Клыкова
\quad$\cdot$\quad
ЕГЭ}
\end{center}

\end{document}